%%%%%%%%%%%%%%%%%%%%%%%%%%%%%%%%%%%%%%%%%%%%%%%%%%%%%%%%
% Este é um documento que servirá de modelo para
% os relatórios feitos na disciplina Laboratório de Circuitos Lógicos
% 2020-2
%%%%%%%%%%%%%%%%%%%%%%%%%%%%%%%%%%%%%%%%%%%%%%%%%%%%%%%%%

%%%%%%%%%%%%%%%%%%%%%%%%%%%%%%%%%%%%%%%%%%%%%%%%%%%%%%%%%
% Use os diferentes diretórios para colocar os relatórios de cada experimento, deste modo vc consegue manter um histórico e todo material organizado em apenas um local.
% Lembre-se de mudar o Main Document no Menu!!!

\documentclass[12pt]{article}

\usepackage{sbc-template}
\usepackage[brazil]{babel}
\usepackage{graphicx}
\usepackage{url}
\usepackage{float}
\usepackage{listings}
\usepackage{color}
%\usepackage{todonotes}
\usepackage{algorithmic}
\usepackage{algorithm}
\usepackage{hyperref}
\hypersetup{
    colorlinks=true,
    linkcolor=black,
    filecolor=magenta,      
    urlcolor=blue,
    citecolor=black
    }
     
\sloppy

\title{Experimento 4\\ 
Circuitos Combinacionais: Comparador de Palavras}

\author{Grupo B2 - Turma T2\\   
        Allan Victor Almeida Faria, 261010546\\
        Caio César Gonçalves Ribeiro,  211038217\\
        Danilo Alves Ribeiro,  242003754\\
}


\address{Dep. Ciência da Computação -- Universidade de Brasília (UnB)\\
  CIC0231 - Laboratório de Circuitos Lógicos\\
  \today
  \email{allanvictor.almeida@gmail.com, caiocesar20032003@gmail.com, daniloalvesjob@gmail.com}
}


\begin{document} 
\maketitle

\selectlanguage{american}

\begin{abstract}

This experiment aimed to develop and simulate digital comparator circuits using the Quartus II software. The circuits were constructed from native logic components and reusable sub-circuits, applying concepts such as modularity, reuse, and encapsulation to the development of electronic systems. The designed circuits were simulated to verify their logical behavior and operation according to their respective truth tables. The results demonstrated the applicability of software-based electronic design tools in the development and verification of digital circuits.

\end{abstract}

\selectlanguage{brazil}

\begin{resumo}

Este experimento teve como objetivo desenvolver e simular circuitos comparadores digitais utilizando o software Quartus II. Os circuitos foram construídos a partir de componentes lógicos nativos e subcircuitos reutilizáveis, aplicando conceitos como modularidade, reuso e encapsulamento ao desenvolvimento de sistemas eletrônicos. Os circuitos projetados foram simulados para verificar seu comportamento lógico e funcionamento de acordo com suas respectivas tabelas verdade. Os resultados demonstraram a aplicabilidade de ferramentas computacionais de projeto eletrônico no desenvolvimento e na verificação de circuitos digitais.

\end{resumo}



\section{Introdução}
\label{sec:Introducao}

O uso de ferramentas de software revolucionou
o desenvolvimento de projetos eletrônicos. Sintetizadores, simuladores
e softwares de esquemáticos permitem a produção e teste de circuitos
em larga escala, além de democratizar o acesso a estudantes e pequenos 
desenvolvedores.

\noindent Nesse sentido, princípios de desenvolvimento de
software, como reuso e encapsulamento são aplicados na prototipação
de circuitos eletrônicos, facilitando e agilizando tal processo. Assim,
foram desenvolvido e simulados no Quartus II circuitos comparadores, utilzando-se
de sub-circuitos reutilizáveis construídos a partir de circuitos nativos.

\subsection{Objetivos}
\label{sec:Objetivos}

Descrever aqui os objetivos do experimento.

\subsection{Materiais} 
\label{sec:Materiais}
Neste experimento foram utilizados os seguintes materiais e equipamentos:
\begin{itemize}
    \item Simulador Quartus II

\end{itemize}

\section{Procedimentos e Resultados}
\label{sec:Procedimentos}

Deve conter a descrição completa de cada item solicitado no experimento, como foram realizados, problemas encontrados e resultados obtidos (esperados ou não). Deve conter figuras, fotos, gráficos, tabelas e links clicáveis para os vídeos solicitados.

\subsection{Multiplexador de 4 entradas}
\label{sec:Mux}



\section{Conclusões}
\label{sec:Conclusao}

Concluir o relatório explanando rapidamente o que foi feito e os resultados obtidos, sempre correlacionando com os objetivos do experimento apresentados na Seção~\ref{sec:Objetivos}. 
\cite{boulic:91}

\bibliographystyle{plain}
\bibliography{relatorio}  %Aqui é a definição do arquivo .bib a ser usado pelas referências


\newpage 
% Colocar aqui apenas as respostas dos itens da Auto-Avaliação
\section*{Auto-Avaliação}

\begin{enumerate}
    \item b
    \item d
    \item d. O circuito é impossível apenas com XOR
    \item d
    \item Se a constante '1' está disponível, "a". Caso contrário, "d" 
\end{enumerate}


\end{document}
